\section{Study Design Explained}

\subsection{Open-label, single-arm, and what that means for you}

\draftinstr{Roughly 1,000 characters. Three terms, each with a consequence.
Open-label means nothing about your treatment is concealed from you. Single-arm
means there is no allocation that could assign you to a comparator, so the
common fear of being randomised to a treatment you would not have chosen does
not arise here. Dose-finding means the study is looking for the highest dose
that is tolerated, not for the largest effect. Cite \texttt{NCITrial2024} for
the general framing and \texttt{onpremwhippl} for the design.}

\subsection{Your status in the study, and how it can change}

\draftinstr{Roughly 900 characters introducing Figure 12 and the idea of a
guard. Every transition in the figure carries a named condition that must
evaluate true, and the protections live in the guards rather than in the states.
Adapt from \texttt{inputs/\pfb phase-\pfb 1-\pfb trial-\pfb protocol.zip},
\texttt{sections/\pfb sec-\pfb 04-\pfb design.tex}.}

\begin{pafig}[plantuml-type: state machine, @startuml]
\node[umlkey] (a) at (0,0) {Figure 12 slot\\participant state machine};
\node[umlbox] (b) at (0,-1.9) {\ttfamily\tiny Source\\ plantuml/\\ fig-\pfb 12-\pfb participant-\pfb state-\pfb machine.puml};
\draw[umlarrow] (a) -- (b);
\end{pafig}
\vspace{-0.7cm}
\figcaption{Figure 12. Participant status as a guarded state machine, screening through the
24-month endpoint, with the sentinel cohort hold drawn as an explicit state because\\
it is the clearest evidence that the design accepts delay in exchange for safety.}

\draftinstr{Replace the Figure 12 slot with the PlantUML-type state machine from
\texttt{plantuml/\pfb fig-\pfb 12-\pfb participant-\pfb state-\pfb machine.puml}. Single column of
\texttt{umlstate} boxes at 2.4 cm pitch; the two off-path states at $x = 6.8$
and $x = -6.8$ aligned with the state they branch from; guards in white-filled
\texttt{umlguard} nodes at the transition midpoint; the self-transition on
\texttt{DoseAssigned} as a 5 mm right-hand loop at looseness 5.}

\subsection{The 3+3 rule, and why you are never the untested dose}

\draftinstr{Roughly 1,300 characters. Explain the rule operationally rather than
statistically. Three participants at a dose level; if none has a dose-limiting
toxicity, the next level opens; if one does, the level expands to six rather
than escalating; if two or more do, escalation stops and the level below becomes
the recommended dose. Add the sentinel stagger: the second participant at a new
level is not dosed until the first participant's safety window has closed. State
that this costs the study time and is not there for the study's benefit.}

\begin{pafig}[graphviz-type: decision tree, record leaves]
\node[gvkey] (a) at (0,0) {Figure 13 slot\\3+3 decision tree};
\node[gvbox] (b) at (0,-2.0) {\ttfamily\tiny Source\\ graphviz/\\ fig-\pfb 13-\pfb escalation-\pfb decision-\pfb tree.dot};
\draw[gvedge] (a) -- (b);
\end{pafig}
\vspace{-0.7cm}
\figcaption{Figure 13. The 3+3 dose-escalation rule drawn as a decision tree with exactly
one path from the root to each leaf, including the two stopping leaves and the study\\
pause that convenes the Data and Safety Monitoring Board and is reported to you.}

\draftinstr{Replace the Figure 13 slot with the graphviz-type decision tree from
\texttt{graphviz/\pfb fig-\pfb 13-\pfb escalation-\pfb decision-\pfb tree.dot}. Four ranks at
$y = 0, -3.0, -6.4, -9.8$; branch labels placed on the upper third of each edge,
not the midpoint, so sibling labels cannot collide; the sentinel-wait annotation
as a \texttt{gvcluster} dashed box behind rank 1; record leaves built from
stacked \texttt{gvcell} rows.}

\subsection{The dose levels}

\draftinstr{Populate the dose levels and the cohort sizes from
\texttt{inputs/\pfb phase-\pfb 1-\pfb trial-\pfb protocol.zip},
\texttt{sections/\pfb sec-\pfb 04-\pfb design.tex} and \texttt{sections/\pfb sec-\pfb 06-\pfb intervention.tex}.
Keep the final column patient-facing.}

\vspace{0.30em}

\noindent\begin{tabularx}{\textwidth}{L{1.9cm} L{3.1cm} C{1.9cm} Y}
\toprule
\textbf{Level} & \textbf{Daraxonrasib dose} & \textbf{Cohort size} & \textbf{What has to be true before this level opens} \\
\midrule
DL1 & Starting dose level & 3, expandable to 6 & Phase 0 gate signed, USL at least 7.0, IND in effect, IDE approved \\
DL2 & One step above DL1 & 3, expandable to 6 & DL1 cleared: 0 of 3, or 1 of 6, dose-limiting toxicities \\
DL3 & Two steps above DL1 & 3, expandable to 6 & DL2 cleared on the same rule, and DSMB review at the cohort boundary \\
\bottomrule
\end{tabularx}

\vspace{0.30em}

\noindent Planned enrollment is up to 18 participants across the three levels.
Enrollment is staggered: no participant enters a level while a sentinel window
at that level remains open.

\subsection{The Phase 0 gate that runs before any patient}

\draftinstr{Roughly 1,100 characters. At least 1000 simulated procedures across
at least two independent simulation frameworks, reproducing motion within a 2 mm
positional and 0.5 N force tolerance, with a documented Unified Safety Level of
at least 7.0. Cite \texttt{asmevv40} and \texttt{nasastd7009} for the credibility
standards this is built against, \texttt{cfr312paiuni} for the Subpart J overlay,
and \texttt{pdac060s2030} for the simulation source. State clearly that a
simulated worst case is not a human worst case.}

\subsection{Stopping rules, stated before they are needed}

\draftinstr{Roughly 900 characters listing the prespecified pause and stopping
triggers from \texttt{inputs/\pfb phase-\pfb 1-\pfb trial-\pfb protocol.zip},
\texttt{sections/\pfb sec-\pfb 04-\pfb design.tex} \S4.6. Include the device-specific triggers,
because these are the ones a participant worried about the robot will look for:
any breach of a force cap, any entry into a vascular no-fly corridor, any
emergency-stop latency above budget, and any loss of the heartbeat bus.}
